\section*{Chapitre \ref{ch:fm:the-asset-manager-and-the-fund} --- La société de gestion et le fonds}

\begin{solution}{exo:fm:the-asset-manager-and-the-fund:1}
Mise en commun~: A détient 100 parts et paie $0.2\times0.05\times100=1.00$~; B détient 111.1 parts et paie 1.11. Séries~: A paie 1.00~; le gain de B est
$100\times(1.05/0.90-1)=16.67$, avec une commission de 3.33.
\end{solution}

\begin{solution}{exo:fm:the-asset-manager-and-the-fund:2}
Investisseurs~: 19.1~\% de la valeur liquidative~; portefeuilles~: 61.1~\%.
\end{solution}

\begin{solution}{exo:fm:the-asset-manager-and-the-fund:3}
2.5~\%, 2.2~\% et 0.5~\% de la valeur liquidative.
\end{solution}

\begin{solution}{exo:fm:the-asset-manager-and-the-fund:4}
$(0.20-0.15)/(1-0.15)=5.9~\%$.
\end{solution}

\begin{solution}{exo:fm:the-asset-manager-and-the-fund:5}
Une commission cristallisée n'est jamais restituée. Sur des mois de $+10\%$ puis $-10\%$~: la \omterm{def:fm:the-asset-manager-and-the-fund:class}{cristallisation} mensuelle prélève $0.2\times0.10=0.02$ du capital dès le
premier mois~; cristallisée sur les deux mois, le fonds finit à 0.99, sous sa valorisation, et ne verse aucune commission.
\end{solution}

\begin{solution}{exo:fm:the-asset-manager-and-the-fund:6}
Le rachat au prorata vend une tranche de chaque compartiment de liquidité dans la semaine, y compris les plus lents, à un coût de vente forcée que supportent les investisseurs restants, faute de \omterm{def:fm:the-asset-manager-and-the-fund:gate}{swing pricing}~;
en contrepartie, le portefeuille restant conserve sa structure et sa liquidité d'origine. La méthode Liquid first ne vend que les actifs faciles à céder et laisse le reste.
\end{solution}

\begin{solution}{exo:fm:the-asset-manager-and-the-fund:7}
Après les premiers 5~\%~: $(8-5)/95=3.16~\%$ vendables en sept jours. Les 5~\% suivants dépassent ce seuil et ne laissent aucune part vendable en sept jours.
\end{solution}

\begin{solution}{exo:fm:the-asset-manager-and-the-fund:8}
La liquidité est une propriété du couple et du chemin~: une première vague de rachats servie selon la méthode Liquid first abaisse le seuil de 60~\% pour la suivante~; les \omterm{def:fm:the-asset-manager-and-the-fund:side}{lettres d'accompagnement} peuvent
accorder à certains investisseurs de meilleures conditions~; le seuil de 60~\% peut être une estimation issue de marchés calmes. Mesurer l'échelle de liquidité après une séquence de rachats sous tension, avec
des coûts de cession en situation de tension.
\end{solution}

\begin{solution}{pb:fm:the-asset-manager-and-the-fund:1}
\begin{enumerate}
\item Les nourriciers, construits pour des investisseurs différents, investissent dans un même maître qui détient le portefeuille~; le \omterm{def:fm:the-asset-manager-and-the-fund:gp}{commandité} gère le fonds et en répond indéfiniment, les
\omterm{def:fm:the-asset-manager-and-the-fund:gp}{commanditaires} apportent leur capital avec une responsabilité limitée.
\item L'\omterm{def:fm:the-asset-manager-and-the-fund:admin}{administrateur de fonds} calcule la valeur liquidative, tient le registre et calcule les commissions~; le \omterm{def:fm:the-asset-manager-and-the-fund:admin}{dépositaire} conserve les actifs. L'indépendance empêche la société de gestion
de valoriser ou d'utiliser elle-même les actifs.
\item Voir les définitions du chapitre.
\item Un \omterm{def:fm:the-asset-manager-and-the-fund:ucits}{OPCVM} (UCITS) rachète à la demande de tout porteur de parts et ne peut suspendre les rachats que temporairement et dans des cas exceptionnels.
\item Investisseurs~: 9.0~\%, portefeuilles~: 51.8~\%.
\item Non~: les portefeuilles sont plus liquides que les conditions de rachat des investisseurs à chaque horizon~; le ratio est le plus faible à 365 jours (76.7 contre 58.2).
\item 2.5~\%, 2.2~\% et 0.5~\%.
\item Seulement 8~\% de ses actifs pouvaient être vendus en sept jours, pour un droit de rachat à quatre jours~; les actifs liquides avaient été vendus de façon disproportionnée.
Les décisions de 2025 sont portées devant le Upper Tribunal et restent provisoires~; le directeur de la société a transigé avec un dispositif d'indemnisation pouvant atteindre 230~millions de livres.
\item Voir la \cref{prop:fm:the-asset-manager-and-the-fund:liquidfirst}.
\item Liquid first~: 0.04~\% et 31.1~\%. Au prorata~: 5.39~\% et 51.8~\%.
\item Liquid first~: 1.69~\% et 0. Au prorata~: 8.81~\% et 8.0~\%.
\item Il règle 10~\%, reporte 20~\%, laissant 46.4~\% (échelle du hedge fund) et 0 (fonds illiquide) vendables en sept jours.
\item Il fait supporter le coût de vente à ceux qui rachètent et supprime la dilution des investisseurs restants~; il ne crée pas de liquidité.
\item 51.8~\% et 8~\% de la valeur liquidative.
\item Des ajustements pour que chaque investisseur ne paie de commission de performance que sur son propre gain~; l'avantage indu de B est de $3.33-1.11=2.22$.
\item 11.54~\%, 11.26~\% et 10.78~\% du capital initial sur cinq ans.
\item L'investisseur bénéficiant d'une \omterm{def:fm:the-asset-manager-and-the-fund:side}{lettre d'accompagnement} peut sortir en premier et laisser les autres avec un fonds moins liquide.
\item Lorsque les investisseurs peuvent demander davantage de rachats sur une période que le fonds ne peut en vendre à un coût normal.
\item Dans les sept jours, à un coût inférieur à 1~\%~: 51.8~\% pour l'échelle Form PF, 8~\% pour le fonds illiquide~; dilution de Liquid first pour une demande de 30~\%
sans \omterm{def:fm:the-asset-manager-and-the-fund:gate}{swing pricing}~: 0.04~\% et 1.69~\% de la valeur liquidative des investisseurs restants.
\item Adapter les conditions aux actifs~: préavis plus long, fréquence de rachat moins élevée ou plafonnement dans la documentation, avant l'arrivée de tout rachat~; et conserver un
coussin de liquidité plutôt que de vendre d'abord les actifs liquides.
\end{enumerate}
\end{solution}

\begin{solution}{iq:fm:the-asset-manager-and-the-fund:1}
Mettre en commun, dans un même portefeuille, des investisseurs à la situation fiscale différente~; le fonds maître détient le portefeuille, et les investisseurs détiennent les nourriciers.
\iqlookfor{la séparation entre le véhicule et la société de gestion.}
\end{solution}

\begin{solution}{iq:fm:the-asset-manager-and-the-fund:2}
Un plafond sur les rachats par date de négociation~; il empêche une ruée de contraindre à des ventes forcées, qui diluerait ceux qui restent.
\iqlookfor{le \omterm{def:fm:the-asset-manager-and-the-fund:gate}{plafonnement des rachats} comme protection des investisseurs restants.}
\end{solution}

\begin{solution}{iq:fm:the-asset-manager-and-the-fund:3}
Avec une seule valeur liquidative par part et une seule valorisation, le nouvel investisseur ne paie aucune commission de performance tant que la valorisation n'est pas regagnée, alors que son propre investissement a progressé~;
la comptabilité par séries ou l'\omterm{def:fm:the-asset-manager-and-the-fund:eq}{égalisation} la lui fait payer sur son propre gain.
\iqlookfor{une seule valorisation pour plusieurs prix d'entrée.}
\end{solution}

\begin{solution}{iq:fm:the-asset-manager-and-the-fund:4}
Il laisse aux investisseurs restants un fonds moins liquide, de sorte que le rachat suivant coûte plus cher et que sortir tôt est récompensé
(voir la \cref{prop:fm:the-asset-manager-and-the-fund:liquidfirst}).
\iqlookfor{l'incitation à la ruée.}
\end{solution}

\begin{solution}{iq:fm:the-asset-manager-and-the-fund:5}
Il fait supporter le coût des flux nets importants aux investisseurs qui les provoquent~; il ne rend pas les actifs plus faciles à vendre.
\iqlookfor{dilution face à la liquidité.}
\end{solution}

\begin{solution}{iq:fm:the-asset-manager-and-the-fund:6}
Comparer, à chaque horizon, la part de la valeur liquidative que les investisseurs peuvent racheter (conditions, plafonnements, \omterm{def:fm:the-asset-manager-and-the-fund:side}{lettres d'accompagnement}) à la part que le portefeuille peut céder en
un délai donné, à un coût annoncé et en situation de tension~; demander l'échelle de liquidité position par position, le registre des investisseurs par conditions et les \omterm{def:fm:the-asset-manager-and-the-fund:side}{lettres d'accompagnement}.
\iqlookfor{les deux côtés du couple, et les scénarios de tension.}
\end{solution}
